\documentclass[12pt,a4paper]{article}
\usepackage[spanish,es-nodecimaldot]{babel}
\usepackage{fontspec}
% Instalar Times New Roman y compilar con XeLaTeX para usar la fuente solicitada.
\setmainfont{Times New Roman}
\usepackage[margin=2.3cm]{geometry}
\usepackage{longtable,array,booktabs,fvextra,enumitem,fancyhdr,needspace}
\usepackage[hidelinks,unicode]{hyperref}
\usepackage{xurl}
\setlength{\parindent}{0pt}
\setlength{\parskip}{6pt}
\setlength{\emergencystretch}{3em}
\setlist{nosep,leftmargin=1.5em,topsep=5pt}
\pagestyle{fancy}\fancyhf{}
\fancyhead[L]{\small IDENTITY ACCESS LAB}\fancyhead[R]{\small Manual de laboratorio}
\fancyfoot[C]{\thepage}
\setlength{\headheight}{15pt}
\DefineVerbatimEnvironment{Comandos}{Verbatim}{breaklines,breakanywhere,fontsize=\small,fontfamily=\rmdefault,frame=single,framesep=5pt}
\begin{document}
{\LARGE\bfseries Dirección y Gestión del Proyecto\par}\vspace{6pt}
2 de octubre de 2026\par\hrule\vspace{10pt}
\Needspace{5\baselineskip}\section*{Objetivo y responsabilidades}
Organizar el trabajo, aprobar reglas del escenario y comprobar que los departamentos entreguen un sistema integrado. Dirección no tiene acceso automático al área restringida ni recibe permisos de administración por su cargo.\par
Puestos: dirección general, responsable de seguridad y coordinación. Cada puesto debe tener una persona asignada en el control privado. Estos cargos no conceden permisos técnicos por sí mismos.\par
\Needspace{5\baselineskip}\section*{Paso 1. Asignar puestos sin publicar datos personales}
\begin{enumerate}
\item Revisar los 31 puestos genéricos de los diagramas.
\item Asignar una persona a cada puesto en un registro privado separado del repositorio.
\item Mantener en privado nombres, correos, tarjetas, aprobadores y contactos.
\item Confirmar que no hay tareas sin responsable ni conteos duplicados.
\item Actualizar los departamentos y funciones genéricas si cambia el diseño.
\item Acordar responsables y suplentes por área en el control privado.
\end{enumerate}
Resultado: 31 puestos con tareas definidas y documentación pública sin datos personales.\par
\Needspace{5\baselineskip}\section*{Paso 2. Registrar decisiones de implementación}
Completar \textnormal{Plantillas/\allowbreak{}acuerdos.csv} con fecha, decisión, responsable y evidencia. Registrar NVR virtual autorizado, alcance del monitoreo móvil, sustitución de transceptores, recursos de Microsoft 365/OneLogin y formato del informe.\par
OneLogin es obligatorio: si el entorno no permite un conector o aprovisionamiento, solicitar a la institución acceso o licencia que lo permita. No aprobar una sustitución por otra plataforma.\par
\Needspace{5\baselineskip}\section*{Paso 3. Asignar equipos}
\begin{enumerate}
\item Confirmar las 14 laptops principales del archivo inicial.
\item Registrar propietario, función, cargador y ubicación.
\item Seleccionar servidor TI y grabador de Seguridad Física. Identificar su sistema operativo y si se preparará Ubuntu en una máquina virtual.
\item Reservar dos teléfonos con sus cargadores y soportes.
\item Confirmar switches, router/punto de acceso, cables y adaptadores.
\item Asignar fecha para identificar físicamente lector y microcontrolador.
\end{enumerate}
Resultado: el inventario indica qué existe y quién lo lleva. “Posiblemente disponible” sigue siendo un pendiente.\par
\Needspace{5\baselineskip}\section*{Paso 4. Crear calendario por resultados}
{\small\renewcommand{\arraystretch}{1.2}
\begin{longtable}{p{7.760cm}p{7.760cm}}
\toprule
\textbf{Etapa} & \textbf{Resultado exigido para cerrar} \\ \midrule\endfirsthead
\textbf{Etapa} & \textbf{Resultado exigido para cerrar} \\ \midrule\endhead
Preparación & Personal, recursos y acceso a plataformas confirmados \\
Infraestructura & Red, aplicación básica y un teléfono accesible \\
Control físico & UID, cuatro tarjetas, LED y CSV funcionando \\
Identidades & OneLogin real, roles, MFA, creación y baja automática \\
Monitoreo & Dos cámaras, móvil, grabación y recuperación de video \\
Integración & Entrada, tarea permitida, tarea denegada, salida y baja \\
Entrega & Pruebas, informe en inglés, presentación y resultados \\
\bottomrule\end{longtable}}
Poner fechas reales en el Excel. Una etapa puede continuar mientras otra espera recursos, pero no declarar completa una dependencia que falla.\par
\Needspace{5\baselineskip}\section*{Paso 5. Aprobar reglas del escenario}
\begin{enumerate}
\item Confirmar las cuatro identidades físicas: LAB-INGENIERO, LAB-ADMIN, LAB-GUARDIA y LAB-AUDITOR.
\item Aprobar funciones digitales: mantenimiento para ingeniero, administración para administrador, consulta/registros para auditor y funciones específicas para Seguridad Física.
\item Establecer que RRHH solicita, el dueño del recurso aprueba y TI/IAM/Seguridad Física configuran.
\item Prohibir la aprobación de un permiso por la misma persona que lo solicita cuando puedan separar funciones.
\item Registrar qué ocurre con una tarjeta desconocida, una baja y un fallo de cámara.
\end{enumerate}
\Needspace{5\baselineskip}\section*{Paso 6. Revisión de progreso}
En cada reunión, preguntar por resultado, evidencia, dificultad y siguiente acción. Actualizar estado y porcentaje en el Excel. Si se declara “Terminado”, abrir su evidencia y comprobarla.\par
Ejemplos de dificultades útiles: “La aplicación recibe un rechazo de callback”; “OneLogin deja el alta en pendiente”; “El iPhone cambia de dirección”. Evitar registros como “no funciona” sin indicar el paso fallido.\par
\Needspace{5\baselineskip}\section*{Paso 7. Coordinar ensayo y exposición}
\begin{enumerate}
\item Comprobar que RRHH dio de alta al ingeniero antes del recorrido.
\item TI, IAM y Seguridad Física verifican el estado inicial.
\item SOC abre imágenes y registros.
\item Mostrar visitante rechazado, ingeniero autorizado, MFA, mantenimiento, auditor rechazado en mantenimiento y salida.
\item Hacer cambio de rol y baja al final para no romper los pasos anteriores.
\item Auditoría presenta resultados reales y limitaciones.
\item Medir duración y ajustar al tiempo del profesor.
\end{enumerate}
Guardar una copia de videos y registros por si falla internet en vivo. Esas evidencias corresponden a un ensayo identificado; no simular que ocurrieron en ese momento.\par
\Needspace{5\baselineskip}\section*{Comandos y archivos}
Dirección trabaja principalmente en el control compartido. Para crear carpetas en una laptop Ubuntu:\par
\Needspace{5\baselineskip}
\begin{Comandos}
mkdir -p ~/lab-runtime/evidencias/direccion
mkdir -p ~/lab-runtime/entrega
\end{Comandos}
En Windows, crear las mismas carpetas con el explorador. No necesitan instalar Docker ni ejecutar comandos administrativos para coordinar el grupo.\par
\Needspace{5\baselineskip}\section*{Entregan}
Lista revisada, decisiones, inventario, calendario, política de acceso, guion y carpeta de entrega. Cada documento indica fecha y responsable.\par
\Needspace{5\baselineskip}\section*{Referencias y videos}
\begin{itemize}
\item \href{https://www.onelogin.com/video/groups-roles-and-mappings-oh-my}{Video oficial OneLogin: grupos, roles y asignación automática}: comprender qué debe pedir Dirección a IAM.
\item \href{https://www.onelogin.com/resource-center/videos/getting-started-with-identity-lifecycle-management-pt-2}{Video oficial OneLogin: ciclo de vida de identidades, parte 2}: identificar resultados de altas y bajas.
\item \href{https://onelogin.service-now.com/kb?id=kb_article_view&sysparm_article=KB0010298}{Introducción oficial al aprovisionamiento}: diferencia entre asignar una aplicación y mantener su cuenta sincronizada.
\item Referencia del proyecto (REFERENCIAS\_Y\_VIDEOS.pdf): fuentes y materiales compartidos.
\end{itemize}
\end{document}
